\section{Лекция 16. Теорема Руше. Геометрические принципы}
\subsection{Теорема Руше}
\textbfit{Теорема} (Руше/ <<Дамы с собачкой>>). Пусть $g, \ f \in \mathcal{O}(D), \ G$ -- ограниченная, $\overline{G} \subset D, \ \partial G$ состоит из одной кусочно-гладкой замкнутой жордановой кривой.
Пусть $|f|>|g|$ на $\partial G$. Тогда $N(f+g, G) = N(f, G)$.
\[\triangleright \ N(f+g, G) = \frac{1}{2\pi} \Delta_{\partial G} \arg (f+g)\]
\[N(f, G) = \frac{1}{2\pi} \Delta_{\partial G} \arg f\]
\[\arg (f+g)(z) = \arg f(z) \cdot \left(1+\frac{g}{f}\right)= \arg f(z) + \arg \left(1+\frac{g}{f}\right) + 2\pi k\]
\[\Delta_{\partial G} \arg (f+g) = \Delta_{\partial G} \arg f + \Delta_{\partial G} \arg \left(1+\frac{g}{f}\right) = N(f, G) \cdot 2\pi + \Delta_{\partial G} \arg \left(1+\frac{g}{f}\right)\]
\[\left|\frac{g}{f}\right| < 1 \Rightarrow 1 + \frac{g}{f}(z) \text{ не обходит 0} \Rightarrow \Delta_{\partial G} \arg \left(1+\frac{g}{f}\right) = 0 \ \blacktriangleleft\]

\textbfit{Теорема} (основная теорема алгебры). Пусть $P_n(z)$ -- многочлен степени $n$ с комплексными коэффициентами. Тогда у него в $ \mathbb{C}$ ровно $n$ нулей с учётом кратности.
\[\triangleright P_n(z) = a_0 z^n + a_1 z^{n-1} + \ldots + a_n\]
\[f = a_0 z^n, \ g = a_1 z^{n-1} + \ldots + a_n\]
\[|g(z)|  \leq n \max_{i=1\ldots n} |a_i| \cdot |z|^{n-1}, \ M = n \max_{i=1\ldots n} |a_i|\]
\[\text{При } R > M \ |f(z)| = |a_0 z^n| > |a_0| \cdot \frac{M}{|a_0|} \cdot |z|^{n-1}, \ \text{т.е. на } |z| = R \ |g| < |f| \Rightarrow\] 
\[\Rightarrow N(P_n(z), |z| < R) = n. \ \blacktriangleleft\]
\subsection{Геометрические принципы}
\textbfit{Теорема} (принцип сохранения области). Пусть $f \in \mathcal{O}(D), \ f \neq const$ в $D, \ D^* = f(D)$. Тогда $D^*$ -- область.
\[\triangleright \text{ 1) Связность. Пусть } A=f(a), \ B = f(b), \ a, b \in D. \text{Т.к. } D \text{ связна } \Rightarrow \exists \text{ пусть } \Gamma:[0, 1] \to D:\]
\[\Gamma(0)=a, \ \Gamma(1) = b.\] 
\[\text{Тогда  } f \circ \Gamma:[0, 1] \to D^*, \ f(\Gamma(0)) = A, \ f(\Gamma(1)) = B \Rightarrow D^* \text{ -- связная область}\]
\[\text{2) Открытость. Пусть }w_0 = f(z_0), \ w_0, f(z_0) \in D^*\]
\[\text{Пусть } \overline{U_\delta (z_0)} \subset D: \ f(z) \neq w_0 \text{ в } \overline{\mathring{U}_\delta}(z_0). \text{ Пусть } r^* = \min_{|z-z_0|=\delta} |f(z) - w_0|.\]
\[\text{Докажем, что } U_{r^*}(w_0) \subset D^*.\]
\[w_1:|w_1-w_0|< r^*\]
\[\text{Хотим } \exists z_1 \in D:f(z_1)=w_1\]
\[f(z) - w_1 = f(z) - w_0 + w_0 - w_1\]
\[\text{На } |z-z_0| = \delta, \ |f(z) - w_0| \geq r^* > |w_0-w_1|, \text{ т.е. } N(f(z)-w_1, |z-z_0|<\delta) =\] 
\[=N(f(z) - w_0, |z-z_0|<\delta) \geq 1 \Rightarrow w_1 \in D^* \ \blacktriangleleft\]

\textbfit{Теорема} (критерий локальной однолистности). Пусть $f \in \mathcal{O}(a).$ Тогда $f$ локально однолистна в $a$ (т.е. $\exists U_\delta(a)$, в которой $f$ -- инъекция) $\Leftrightarrow f'(a) \neq 0.$
\[\triangleright \ (\Leftarrow) \text{ уже доказывали (т. об обратном отображении + Коши-Римана)}\]
\[(\Rightarrow) \text{ Пусть $f$ -- инъекция в } U(a), \ f'(a) = 0\]
\[f(a) = w_0. \ f(z) - w_0 \text{ имеет в $a$ 0 порядка $p \geq 2$.}\]
\[\text{По т. о ! можем выбрать } U_\epsilon(a): \ f'(z) \neq 0 \text{ в } \mathring{U}_\epsilon(a), \ f(z) \neq w_0.\]
\[\text{Рассмотрим } U_\delta(w_0) \subset D^* = f(U_\epsilon(a)), \text{ возьмём } w_1 \in U_\delta(w_0).\]
\[r^* := \min_{|z-a|=\epsilon} |f(z) - w_0|, \ \text{пусть } w_1 \in U_\delta(w_0) : \ |w_1-w_0|<r^*.\]
\[\text{Тогда } f(z) - w_1 = f(z) - w_0 + w_0 - w_1\]
\[\text{На } |z-a|=\epsilon \ |f(z) - w_0|\geq r^* > |w_0-w_1| \Rightarrow N(f(z)-w_1, |z-a|<\epsilon) =\] 
\[=N(f(z)-w_0, |z-a|<\epsilon) = p \geq 2.\]
\[\text{Т.к. } f'(z)\neq 0 \text{ в } U_\epsilon(a), \text{ то все нули $f(z) - w_1$ простые} \Rightarrow \text{ их в точности $p$. } \blacktriangleleft\]

\textbfit{Теорема} (локальный принцип максимума модуля). Пусть $f \in \mathcal{O}(D)$ и $a \in D$ -- точка локального максимума $|f|$. Тогда $f \equiv const$ в $D$.
\[\triangleright \text{ Пусть } \forall z \in U_\epsilon(a) \ |f(z)|<|f(a)|. \]
\[D^* = f(U_\epsilon(a)) \text{ -- область. } w_0 = f(a)\]
\[\text{Тогда } \exists U_\delta(w_0) \subset D^*\text{, но } \exists w_1 \in U_\delta(w_0) : \ |w_1|>|w_0|\]
\[w_1 \in D^* \Rightarrow w_1 = f(z_1), \ z_1 \in U_\epsilon(a) \text{ -- противоречие } \blacktriangleleft\]

\textbfit{Теорема} (принцип максимума модуля для ограниченной функции). Пусть $f \in \mathcal{O}(D), \ f \in C(\overline{D}), \ D$ -- ограниченная область. Тогда $\max_{\overline{D}} |f(z)| = \max_{\partial D} |f(z)|.$
\[\triangleright \ \text{Если $f \equiv const$, то выполнено. Иначе достигаться во внутренней точке не может, но на}\]
\[\overline{D} \text{ достигается (Вейерштрасс)} \Rightarrow \text{ на } \partial D \blacktriangleleft\]

\textbfit{Утв.} (принцип минимума). Пусть $f \in \mathcal{O}(D), \ f \in C(\overline{D}), \ D$  ограничена, то или $\exists z_0 \in D: |f(z_0)| = 0 (\min_D |f(z)| = 0)$, или $\min_{\overline{D}} |f(z)|  = \min_{\partial D} |f(z)|.$
\[\triangleright \text{Если } f(z) \neq 0, \text{ то рассмотрим } \frac{1}{f} \blacktriangleleft\]

\textbfit{Пр.} $\sin z$ в верхней полуплоскости. Максимум всей области $\infty$, но на границе только 1.

\textbfit{Теорема} (Римана). Пусть $D$ -- односвязная область в $\mathbb{C}, \ D \neq \mathbb{C}$. Тогда $\exists f(z) : D \to \mathbb{C},$ конформное в $D: f(D) = \{|z|<1\}$.

$\mathbb{C}$ нельзя, потому что Лиувиль.